\section{Experimental Design}
\label{sec:experimental}

\subsection{Evaluation objectives}

The evaluation is designed to examine the complete CAGE-S decision
pipeline rather than only the predictive model. In particular, the
experiments address the questions summarized in
Table~\ref{tab:research-questions}.

\begin{table*}[t]
\centering
\caption{Research questions guiding the experimental evaluation.}
\label{tab:research-questions}
\begin{tabularx}{\textwidth}{lX}
\toprule
Question & Evaluation objective \\
\midrule
RQ1 &
Can recent temporal telemetry predict next-interval node risk more
accurately and more reliably than current-state telemetry alone? \\

RQ2 &
How does the complete scheduler compare with random, least-loaded,
deterministic-guardrail, current-state learning, policy-only, and
component-ablated alternatives under an identical replay protocol? \\

RQ3 &
Does fairness-aware shortlisting control destination concentration
without changing the operational quality of the selected placements? \\

RQ4 &
How robust are the decisions to service-frequency imbalance, policy
assignment, shortlist size, risk operating point, and destination
reservation capacity? \\

RQ5 &
Can the complete feature-construction, inference, filtering, ranking,
and reservation path execute within an online control interval for a
large candidate set? \\
\bottomrule
\end{tabularx}
\end{table*}

\subsection{Source trace and integrity checks}

The evaluation uses the public Alibaba Microservices Trace 2021
\citep{luo2021alibaba}. The trace contains node-level telemetry and
microservice-instance resource signals sampled at 30-second intervals.
We process the complete \texttt{Node\_0} archive together with the
\texttt{MSResource\_0} shard.

The raw node archive contains 18,868,400 records, 13,118 nodes, and
1,440 timestamps covering a 12-hour observation period. Each
node--timestamp pair is unique. The raw archive contains 2,882 missing
CPU values and 1,440 missing memory values. Rather than interpolating
incomplete node trajectories, strict cleaning retains only nodes with a
complete sequence of CPU and memory observations. The resulting strict
node trace contains 18,803,520 records for 13,058 nodes, with no missing
values or duplicate node--timestamp pairs. The raw microservice shard contains 11,564,394 records, 1,303 services,
96,380 service--instance pairs, 12,731 nodes, and 120 timestamps. A
total of 1,080 records with missing instance CPU or memory originate
from nine service instances. These instances are removed. The remaining
records are aligned with nodes that are valid within the local
microservice-shard interval. The clean microservice trace contains
11,561,072 records, 96,346 service--instance pairs, and 12,727 nodes.
Twenty-two service--instance pairs that enter after the first timestamp
are retained from their first valid observation onward. The node trace is then restricted to the 12,727 nodes and 120 timestamps
represented in the clean microservice shard. This produces an aligned
node table of 1,527,240 records, equal to
$12{,}727\times120$, with no missing observations or duplicate keys.
No CPU or memory value is clipped or imputed during preprocessing.

Table~\ref{tab:data-preparation} summarizes the resulting datasets.

\begin{table*}[t]
\centering
\caption{Trace preparation and replay datasets.}
\label{tab:data-preparation}
\begin{adjustbox}{max width=\textwidth}
\begin{tabular}{lrrrrr}
\toprule
Dataset &
Records/events &
Nodes &
Services &
Service--instance pairs &
Timestamps \\
\midrule
Raw \texttt{Node\_0} &
18,868,400 &
13,118 &
-- &
-- &
1,440 \\

Strict node trace &
18,803,520 &
13,058 &
-- &
-- &
1,440 \\

Raw \texttt{MSResource\_0} &
11,564,394 &
12,731 &
1,303 &
96,380 &
120 \\

Clean aligned \texttt{MSResource\_0} &
11,561,072 &
12,727 &
1,303 &
96,346 &
120 \\

Aligned node snapshots &
1,527,240 &
12,727 &
-- &
-- &
120 \\

Natural evaluable event pool &
20,324 &
12,727 &
368 &
-- &
108 \\

Service-balanced evaluable pool &
6,875 &
12,727 &
294/285\tnote{a} &
-- &
108 \\
\bottomrule
\end{tabular}
\end{adjustbox}

\begin{tablenotes}
\footnotesize
\item[a] Number of represented services in the calibration/test
partitions, respectively.
\end{tablenotes}
\end{table*}

Instance-level CPU and memory values are used only to identify changes
in service-instance activity. Node admission, temporal prediction, and
guardrail computation use the corresponding node-level telemetry.
This separation ensures that instance signals are not interpreted as an
absolute accounting of node capacity.

\subsection{Replay-event construction}

Almost all service instances are already active at the beginning of the
microservice shard: 96,341 of the 96,346 clean service--instance pairs
appear at the first timestamp. Native service arrivals are therefore too
rare to provide a sufficiently large online scheduling workload. We construct decision events from positive instance-level CPU bursts.
Let $v^{c}_{i,t}$ denote the CPU signal of service instance $i$ at
timestamp $t$. The one-step change is

\begin{equation}
\label{eq:event-cpu-change}
    \Delta v^{c}_{i,t}
    =
    v^{c}_{i,t}
    -
    v^{c}_{i,t-\Delta},
\end{equation}

where $\Delta=30$ seconds. After an initial 300-second warm-up, an event
is emitted when

\begin{equation}
\label{eq:event-rule}
    \Delta v^{c}_{i,t}>0.05.
\end{equation}

The event threshold is applied identically to every service instance and
is fixed before model fitting and scheduler comparison. The warm-up
ensures that temporal histories are available before events enter the
replay.

The initial rule produces 20,419 events over 109 timestamps. Events at
the final timestamp cannot be paired with a subsequent node observation
at $t+\Delta$ and are therefore excluded from evaluable replay. The
final natural event pool contains 20,324 events over 108 timestamps,
representing 368 services and 4,939 observed source nodes.

For each event, the stored replay record contains the service,
service-instance identifier, source node, timestamp, previous and current
instance signals, the instance-level CPU change, and the source-node
state available at the decision timestamp.

\subsection{Temporal alignment and information boundary}

Figure~\ref{fig:timeline} illustrates the temporal protocol. An event is
identified from the service-instance change ending at timestamp $t$.
The scheduler observes the node snapshot at $t$ and temporal features
constructed only from $t$, $t-\Delta$, and $t-2\Delta$. The selected
node's recorded state at $t+\Delta$ is used to compute the native
next-state evaluation metrics.

\begin{figure*}[t]
  \centering
  \includegraphics[width=0.88\textwidth]{temporal_protocol.pdf}
  \caption{Temporally ordered replay protocol. Event detection ends at
  timestamp $t$; scheduling uses only the current and preceding
  observations; evaluation uses the selected node state at
  $t+\Delta$, where $\Delta=30$ seconds.}
  \label{fig:timeline}
\end{figure*}

The following information boundary is enforced throughout all model and
scheduler experiments:

\begin{equation}
\label{eq:information-boundary}
    \mathcal{I}_{t}
    =
    \left\{
        X_{\tau},
        \text{event data}_{\tau}
        :
        \tau\leq t
    \right\}.
\end{equation}

No feature available to a decision at $t$ is constructed from
$t+\Delta$ or any later timestamp. Training labels, calibration
residuals, probability thresholds, shortlist settings, and guardrail
margins are fixed without using the final test region.

\subsection{Natural and service-balanced workloads}

The natural event pool preserves the service frequency distribution
produced by the burst rule. Its event frequency is concentrated: the
ten most frequent services account for approximately 46\% of events.
This workload is used as the primary evaluation because it retains the
observed service imbalance.

To verify that the conclusions are not dominated by a small number of
high-frequency services, we construct a service-balanced workload.
Within each service and timestamp, at most two events are retained:

\begin{equation}
\label{eq:balanced-cap}
    \left|
        \left\{
        e\in\Eset_t:s_e=s
        \right\}
    \right|
    \leq2.
\end{equation}

Selection under this rule is deterministic. The resulting balanced pool
contains 6,875 evaluable events. Its test partition contains 3,787 events,
and the top-ten service share decreases to 22.55\%. The value two in this
workload-construction rule is unrelated to the destination reservation
capacity $C$ used by the scheduler.

\subsection{Chronological development, validation, and test partitions}

The replay interval is split chronologically at timestamp

\begin{equation}
    t_{\mathrm{split}}=1{,}800{,}000~\mathrm{ms}.
\end{equation}

Node transitions before the split form the pre-test region, and
transitions after the split form the final test region. The aligned
temporal-risk dataset contains 1,374,516 node transitions:

\begin{table}[t]
\centering
\caption{Chronological partitions used for model development and final
evaluation.}
\label{tab:partitions}
\begin{tabular}{lrr}
\toprule
Partition & Node transitions & Replay events \\
\midrule
Development & 432,718 & -- \\
Validation & 190,905 & -- \\
Calibration total & 623,623 & 8,695 \\
Final test & 750,893 & 11,629 \\
\bottomrule
\end{tabular}
\end{table}

The development and validation partitions are contained entirely within
the pre-test region. Model fitting is performed on development data.
Estimator selection and operating-point selection use validation or
calibration data. The final test transitions and test replay events are
accessed only after the model, thresholds, margins, and scheduling
configuration are frozen.

For the service-balanced workload, the corresponding calibration and
test event counts are 3,088 and 3,787, respectively. Both natural and
balanced test workloads contain the same 59 decision timestamps.

\subsection{Temporal-risk model development}

The primary risk estimator is a histogram gradient-boosting classifier
implemented with scikit-learn \citep{pedregosa2011sklearn}. Its input is
the 22-dimensional temporal feature vector defined in
Section~\ref{sec:method} and Appendix~\ref{app:features}. The exact
feature order is fixed for training, serialization, and online inference.

A current-state histogram gradient-boosting model is trained as the
direct predictive ablation. It receives current node telemetry and
current-state resource features but excludes temporal lags, differences,
ranges, and variation. Both models use the same development,
validation, and final-test partitions.

Model selection is based on complementary discrimination and probability
quality measures: ROC-AUC, PR-AUC, Brier score, log loss, and expected
calibration error. Scheduler-level thresholds are not selected by
maximizing test performance. Instead, thresholds are determined from
the calibration region at matched operating points.

The default temporal-risk threshold is

\begin{equation}
\label{eq:default-temporal-threshold}
    \theta_{\mathrm{temp}}
    =
    0.011362695354780319.
\end{equation}

The matched current-state HGB threshold is

\begin{equation}
\label{eq:current-hgb-threshold}
    \theta_{\mathrm{current}}
    =
    0.004292615073088162.
\end{equation}

A stricter temporal operating point,

\begin{equation}
    \theta_{\mathrm{safety}}
    =
    0.0020594307448572016,
\end{equation}

is retained for sensitivity analysis. Once selected, these thresholds
remain fixed across workloads and policy seeds.

\subsection{Q95 guardrail calibration}

The persistence baseline predicts that the next node state equals the
current state. For each calibration transition, CPU and memory residuals
are computed as

\begin{align}
    r^{c}_{n,t}
    &=
    \left|
    c_{n,t+\Delta}-c_{n,t}
    \right|,\\
    r^{m}_{n,t}
    &=
    \left|
    m_{n,t+\Delta}-m_{n,t}
    \right|.
\end{align}

The 95th-percentile residual margins computed exclusively from calibration
transitions are

\begin{equation}
\label{eq:q95-values}
    q_c^{0.95}
    =
    0.025929729975254,
    \qquad
    q_m^{0.95}
    =
    0.00021962105567108203.
\end{equation}

These fixed margins are used to construct the CPU and memory upper
projections in Eq.~\eqref{eq:guardrail}. The default operational
boundaries are

\begin{equation}
    \rho_c=\rho_m=0.8.
\end{equation}

The principal guardrail accepts a node when $g_{n,t}\leq1$. Additional
cutoffs of $0.99$ and $0.975$ are evaluated as deterministic sensitivity
profiles.

\subsection{Policy-overlay construction}
\label{sec:policyoverlay}

The trace contains operational telemetry but no jurisdiction, trust,
encryption, trusted-execution, or service-sensitivity attributes.
CAGE-S therefore uses a deterministic policy overlay to evaluate the
policy semantics of the scheduler.

Node identifiers are mapped through a stable cryptographic hash to one
of four jurisdictions---EU-West, EU-Central, US-East, and
AP-Southeast---one of three ordered trust levels, and Boolean encryption
and trusted-execution capabilities. Service identifiers are independently
mapped to one of six policy classes. Hashing depends only on the
identifier and policy seed; it does not use resource telemetry, event
frequency, source placement, model predictions, or future node states.

Table~\ref{tab:policyclasses} reports the number of legal nodes for each
class under the reference seed.

\begin{table}[t]
\centering
\caption{Policy classes and legal-node counts under the reference policy
seed.}
\label{tab:policyclasses}
\begin{tabular}{lr}
\toprule
Service policy class & Legal nodes \\
\midrule
Global public & 12,727 \\
Global protected & 10,320 \\
EU restricted & 6,206 \\
US restricted & 2,613 \\
APAC restricted & 1,501 \\
EU confidential & 866 \\
\bottomrule
\end{tabular}
\end{table}

The complete replay is repeated over 30 deterministic policy seeds.
Changing the seed modifies both service requirements and node
capabilities while leaving the event sequence, telemetry, temporal model,
thresholds, guardrail margins, and scheduling implementation unchanged.
This design isolates sensitivity to policy composition.

\subsection{Compared scheduling methods}

All methods receive identical event batches, timestamp ordering, node
snapshots, active-source sets, destination-capacity rules, and fixed
tie-breaking identifiers. Methods that use policy information receive
the same policy overlay. Except for the observed-source baseline, each
method first retains a source that satisfies its own admission rule and
invokes destination selection only when that source fails.

Table~\ref{tab:compared-methods} defines the complete comparison set.

\begin{table*}[t]
\centering
\caption{Compared methods and enabled scheduling components.}
\label{tab:compared-methods}
\begin{adjustbox}{max width=\textwidth}
\begin{tabular}{lcccccl}
\toprule
Method &
Policy &
Temporal model &
Q95 check &
Fairness &
Reservation &
Destination rule \\
\midrule
Observed source &
\no & \no & \no & \no & \no &
Retain native source \\

Random-policy-Q95 &
\yes & \no & \yes & \no & \yes &
Uniform random accepted destination \\

Least-loaded-policy-Q95 &
\yes & \no & \yes & \no & \yes &
Minimum current utilization \\

Policy-only-fair &
\yes & \no & \no & \yes & \yes &
Fair shortlist; current-state ranking \\

Policy-Q95-fair &
\yes & \no & \yes & \yes & \yes &
Fair shortlist; guardrail ranking \\

Policy-boundary-0.99 &
\yes & \no & \yes & \yes & \yes &
Guardrail cutoff $0.99$ \\

Policy-current-HGB-Q95 &
\yes & Current only & \yes & \yes & \yes &
Current-state learned risk \\

Temporal-Q95-no-policy &
\no & Temporal & \yes & \yes & \yes &
Temporal risk without policy gating \\

Policy-temporal-no-projection &
\yes & Temporal & \no & \yes & \yes &
CAGE-S without final Q95 projection \\

CAGE-S &
\yes & Temporal & \yes & \yes & \yes &
Complete lexicographic pipeline \\
\bottomrule
\end{tabular}
\end{adjustbox}
\end{table*}

The random baseline is evaluated over 20 deterministic replicates. Its
random seed is derived from the workload and replicate identifier, making
every run reproducible. All other methods are deterministic.

\subsection{Parameter selection and sensitivity matrix}

Table~\ref{tab:experimental-parameters} lists the principal experimental
parameters, their default values, and the data used for their selection.
No parameter is tuned on the final test workload.

\begin{table*}[t]
\centering
\caption{Experimental parameters and selection protocol.}
\label{tab:experimental-parameters}
\begin{tabularx}{\textwidth}{lclX}
\toprule
Parameter & Symbol & Value(s) & Selection source \\
\midrule
Control interval &
$\Delta$ &
30~s &
Native sampling interval of the trace \\

Warm-up duration &
-- &
300~s &
Ensures availability of temporal history \\

Event CPU-change threshold &
-- &
$0.05$ &
Fixed burst rule used across all splits and methods \\

CPU/memory boundary &
$\rho_c,\rho_m$ &
$0.8,0.8$ &
Operational headroom boundary \\

Temporal threshold &
$\theta_{\mathrm{temp}}$ &
$0.011362695354780319$ &
Calibration-side matched-availability operating point \\

Current-HGB threshold &
$\theta_{\mathrm{current}}$ &
$0.004292615073088162$ &
Matched operating point for direct model comparison \\

Strict temporal threshold &
$\theta_{\mathrm{safety}}$ &
$0.0020594307448572016$ &
Calibration-side safety sensitivity profile \\

CPU Q95 margin &
$q_c^{0.95}$ &
$0.025929729975254$ &
95th percentile of calibration residuals \\

Memory Q95 margin &
$q_m^{0.95}$ &
$0.00021962105567108203$ &
95th percentile of calibration residuals \\

Shortlist size &
$k$ &
256 &
Validation comparison with
$\{128,256,512,1024,2048,\text{full}\}$ \\

Destination capacity &
$C$ &
1 &
Observed event-concurrency audit; $C=2$ used for sensitivity \\

Policy seeds &
-- &
30 &
Overlay-composition robustness \\

Random replicates &
-- &
20 &
Random-baseline variability \\

Bootstrap replicates &
-- &
20,000 &
Timestamp-cluster uncertainty estimation \\
\bottomrule
\end{tabularx}
\end{table*}

Sensitivity analysis varies one controlled factor at a time while holding
the event sequence and all remaining settings fixed. The examined factors
are workload composition, temporal operating point, deterministic
guardrail cutoff, shortlist size, fairness-aware versus score-only
ranking, destination capacity, and policy seed.

\subsection{Evaluation metrics}

All event-level metrics are computed from the final selected node,
including retained sources and migrated destinations. Let
$n^{*}_{e}$ denote the selected node for event $e$, and let $n_e$ denote
its observed source. Table~\ref{tab:metrics} defines the principal
metrics.

\begin{table*}[t]
\centering
\caption{Evaluation metrics. Lower is better unless stated otherwise.}
\label{tab:metrics}
\begin{tabularx}{\textwidth}{lXc}
\toprule
Metric & Definition & Preferred direction \\
\midrule
Migration rate &
$\frac{1}{|\Eset|}
\sum_{e\in\Eset}
\indicator[n^{*}_{e}\neq n_e]$ &
Lower \\

Policy-violation rate &
$\frac{1}{|\Eset|}
\sum_{e\in\Eset}
\indicator[L(s_e,n^{*}_{e})=0]$ &
Lower \\

Temporal-threshold violation &
$\frac{1}{|\Eset|}
\sum_{e\in\Eset}
\indicator[p_{n^{*}_{e},t_e}>\theta]$ &
Lower \\

Q95-violation rate &
$\frac{1}{|\Eset|}
\sum_{e\in\Eset}
\indicator[g_{n^{*}_{e},t_e}>1]$ &
Lower \\

Native next-state unsafe rate &
$\frac{1}{|\Eset|}
\sum_{e\in\Eset}
y_{n^{*}_{e},t_e+\Delta}$ &
Lower \\

Fallback rate &
Fraction of events that require expansion beyond the initial shortlist &
Lower \\

Infeasible-event rate &
Fraction of events for which no admissible destination can be committed &
Lower \\

Unique destinations &
Number of distinct selected destination nodes &
Higher \\

Top-ten destination share &
Fraction of migrated assignments received by the ten most-used
destinations &
Lower \\

Jain destination fairness &
Distributional fairness of migrated assignments across all $M$ nodes &
Higher \\

Decision latency &
Elapsed online-core processing time per timestamp, reported in ms &
Lower \\
\bottomrule
\end{tabularx}
\end{table*}

Let $z_n$ denote the number of migrated events assigned to node $n$.
Jain's fairness index is computed over the full node set, including nodes
with zero assignments:

\begin{equation}
\label{eq:jain}
    J
    =
    \frac{
        \left(
        \sum_{n=1}^{M}z_n
        \right)^2
    }{
        M
        \sum_{n=1}^{M}z_n^2
    }.
\end{equation}

This definition penalizes concentration on a small subset of the
available infrastructure. In addition to Jain fairness, we report the
top-one and top-ten destination shares and the number of unique
destinations.

\subsection{Statistical analysis}

Events at the same timestamp share a node snapshot, candidate state, and
reservation context and are therefore not treated as independent
observations. Uncertainty for event-level comparisons is estimated using
a paired timestamp-cluster bootstrap. For each workload, the 59 test timestamps are sampled with replacement.
Whenever a timestamp is selected, all events associated with that
timestamp are retained together. The same sampled timestamps are applied
to CAGE-S and the corresponding baseline, preserving pairing. For the
random baseline, replicate identity is additionally sampled inside each
bootstrap iteration.

We generate 20,000 bootstrap samples. For each metric, the reported
effect is

\begin{equation}
\label{eq:bootstrap-difference}
    \Delta_{\mathrm{metric}}
    =
    \mathrm{metric}_{\mathrm{CAGE\text{-}S}}
    -
    \mathrm{metric}_{\mathrm{baseline}}.
\end{equation}

Percentile 95\% confidence intervals are obtained from the
$2.5$th and $97.5$th percentiles of the paired bootstrap distribution.
For metrics where lower values are preferred, a confidence interval
entirely below zero supports a lower value for CAGE-S. For fairness
metrics, the interpretation is reversed.

Policy-seed robustness is summarized separately across the 30
deterministic overlays using the mean, median, and 5th and 95th
percentiles. Random-baseline variability is summarized across its 20
replicates.

\subsection{Runtime measurement}

The online-core benchmark measures the complete decision computation
after the required node snapshots are available in memory. The measured
path includes:

\begin{equation}
\begin{split}
\text{online core}={}&
\text{22-feature construction}
+\text{source inference}\\
&+\text{policy filtering}
+\text{fair shortlisting}\\
&+\text{destination inference}
+\text{ranking}\\
&+\text{reservation}
+\text{pre-commit verification}.
\end{split}
\end{equation}

Risk inference is performed on the timestamp-level superset required by
the event batch. Consequently, the benchmark includes both feature
construction and all model calls needed for the actual decisions rather
than timing a single isolated prediction.

A latency observation is recorded for every test timestamp. We report
the mean, median, 95th percentile, 99th percentile, and maximum elapsed
time. Telemetry acquisition, file and network I/O, serialization, and
external orchestration actions are outside the measured online core.

\subsection{Implementation and reproducibility}

Preprocessing and replay are implemented in Python 3.12.3. The principal
software versions are reported in Table~\ref{tab:software-environment}.

\begin{table}[t]
\centering
\caption{Software environment.}
\label{tab:software-environment}
\begin{tabular}{ll}
\toprule
Component & Version \\
\midrule
Operating system & Ubuntu 24.04.4 LTS under WSL2 \\
WSL2 kernel & 6.18.33.2-microsoft-standard-WSL2 \\
Python & 3.12.3 \\
Polars & 1.42.1 \\
PyArrow & 25.0.0 \\
DuckDB & 1.5.4 \\
pandas & 3.0.3 \\
NumPy & 2.5.1 \\
SciPy & 1.18.0 \\
scikit-learn & 1.9.0 \\
\bottomrule
\end{tabular}
\end{table}

Experiments are executed on a Dell Precision 3561 workstation. The
feature order, model artifact, policy seeds, random seeds, event ordering,
thresholds, and candidate tie-breaking rule are fixed. Intermediate
tables are stored as sorted Parquet files with explicit schemas, and
principal generated artifacts are accompanied by SHA-256 checksums.
All reported scheduler comparisons are produced from the same immutable
test event files and node snapshots.